%% Supplementary material for "When Query Rewriting Loses the Documents".
%% Protocol and comparability detail moved out of the main text; nothing here
%% is new, and every number it explains is reported in the paper itself.
\documentclass[acmsmall,anonymous,screen]{acmart}
\settopmatter{printacmref=false}
\renewcommand\footnotetextcopyrightpermission[1]{}
\pagestyle{plain}
\usepackage{booktabs}
\usepackage{amsmath}
\begin{document}
\title{Supplementary Material: Measurement Protocol and Arm Comparability}
\author{Anonymous Author(s)}
\maketitle

\noindent This supplement holds the protocol detail for \emph{When Query
Rewriting Loses the Documents: An Adaptive Multi-Agent RAG Architecture for
Software Ecosystem Update Monitoring}. Section~S1 traces one question from the
benchmark to one number in a results table. Section~S2 states what makes two
arms comparable. Both are referenced from the main text and neither introduces
a measurement that is not reported there.

\section*{S1\quad How a Reported Number Is Produced}
Every retrieval number in Section~\ref{sec:results} is the mean over questions
of a per-question, per-arm quantity computed by one harness
(\texttt{eval\_harness/run\_eval.py}) in one pass. Because the arms differ in
what they fetch, and because an LLM supplies the relevance labels, the order in
which those steps happen decides what the number can be read to mean. We
therefore trace a single question end to end rather than name the metrics and
leave the protocol implicit.

\paragraph{1. The question.} A benchmark record carries an id, the question
text, its ecosystem and category, the endpoint and record type it was mined
from, and---for a community-mined question---the URL of the post whose title it
is. Every arm in the run receives the same question string.

\paragraph{2. Retrieval.} Each arm retrieves through the same
\texttt{RetrieverAgent} against the endpoints of
Section~\ref{sec:datasources}: the baseline on the question, the multi-agent
arm on the union of the question and its rewriting
(Section~\ref{sec:union-fetch}). Before any ranking, the record's own source
URL is struck from the candidate pool of \emph{every} arm
(Section~\ref{sec:leak}). A run reads its documents in one of two corpus
modes, and which one decides what its numbers can later be re-derived from. In
\emph{record} mode the endpoints are queried live and every response is written
to a snapshot directory; in \emph{replay} mode the fetch is served from that
directory, and in the \emph{strict} variant a document request the snapshot
does not hold raises rather than quietly going live. The raise was not the
guarantee it first appears to be, and we correct our own earlier description of
it. In the configuration every run reported here was made under, each source
fetch in the Retriever is individually guarded, so the exception was caught
there and that source returned no documents for that question: a strict run did
not stop on a miss, it lost a source quietly and reported the count at the end.
That count is where \texttt{frozen} in a run's configuration comes from, and it
is the field to read --- a run whose \texttt{frozen} is false has holes of this
kind, and their number is the honest measure of how frozen it was. We state the
count wherever it is not zero rather than let the mode's name stand in for the
property. Either way the
arms of a single run are compared against one corpus, which is what every
comparison in this paper rests on. Only replay makes a run reproducible
afterwards, and the asymmetry is worth stating precisely because it is easy to
assume otherwise: a snapshot holds one response body per request key, a live
feed answers the same key differently over the hours a 500-question run takes,
and the directory keeps the last of those answers. A record-mode run therefore
cannot be replayed back into itself. Replaying one of ours reproduced its
multi-agent arm's candidate pools on 1 of 12 questions, while two replays of a
single configuration agree on 12 of 12 and are byte-identical. Record-mode
numbers are internally valid and externally irreproducible; replay-mode numbers
are both. Where a table's arms come from different runs we say so, and we do
not read a difference between a live-recorded arm and a replayed one as an
effect of the architecture.

Retrieval also reads the clock, which a snapshot of HTTP responses does not
freeze: a question carrying \emph{latest} or \emph{last week} becomes a date
filter, so the same snapshot replayed a day later can admit a different set of
documents. A recording now stamps its directory with the date it was made and
a replay pins the process clock to that stamp, which removes the drift between
two runs made on different days; an explicitly set clock overrides it, and a
snapshot recorded before the stamp existed reports no recording date rather
than a guessed one.

\paragraph{3. Ranking and cut.} The candidate pool is ordered by the run's
ranking function and cut to $k=4$. Two settings decide the ordering and are
recorded per run: the scorer (BM25, embedding cosine, reciprocal-rank fusion,
or an LLM grading cascade) and whether the verified-source tier prior is
applied before it or not (Section~\ref{sec:rerank}). The top-$k$ list of
document ids is what the metrics see; the documents themselves are what the
answer is synthesized from.

\paragraph{4. Judging.} The judged pool for a question is the union of what
\emph{all} arms returned for it---optionally widened to everything they
considered, which is what makes pool recall computable. Each
(question,~document) pair is labelled once, by \texttt{llama3.1} at temperature
zero, on a graded scale: 2 highly relevant, 1 related, 0 irrelevant. The label
is keyed by a hash of the question text together with the document id and
cached on disk, so it is computed once and reused for every arm in the run and
for every later run that retrieves the same document for the same question. No
arm is scored against a private set of judgments, and no arm is penalized for a
document that only a competing arm surfaced.

\paragraph{5. Metrics.} From the ranked list and those labels:
nDCG@$k$ with gain $2^g-1$ and the standard $\log_2$ discount, its ideal
computed from the same per-question pool; Recall@$k$ over the relevant
documents in that pool; MRR on the first document with grade $\geq 1$. The
denominators are therefore \emph{pool-relative}. A system cannot be charged for
relevant documents no arm retrieved, which makes these numbers valid for the
paired within-run comparisons this paper draws and invalid as absolute values
to compare against published figures computed over a fixed judged corpus. We
state the pooling protocol with every absolute number for this reason, and
Section~\ref{sec:pool-demo} measures what ignoring it costs: two arms whose
retrieved lists are byte-identical, scored in pools of different width, differ
by $+0.055$ recall@5.

\paragraph{6. Answer scoring.} The same judge scores each arm's answer for
faithfulness to its own retrieved contexts and for answer relevance, each in
$[0,1]$, and for correctness against the reference answer on the questions that
carry one. Section~\ref{sec:metric-threat} states what an LLM judge scoring a
template against prose measures, and why the \texttt{marag\_llm} arm exists.

\paragraph{7. Cost.} The harness counts, per arm and per question, every model
call and every prompt and completion token the arm itself spent, attributed by
role---rewriting, ranking, synthesis---with the judge's own calls excluded, so
an accuracy difference between arms can be read against what it cost to obtain.

\paragraph{8. Statistics.} Arms are compared pairwise on the per-question
values, never on the difference of two means: the paired difference vector
feeds an exact two-sided Wilcoxon signed-rank test at the sample sizes where it
is exact, a percentile bootstrap 95\% confidence interval over question-level
differences (10{,}000 resamples, fixed seed), and a win/tie/loss count.
Families of comparisons carry a Holm correction.

\paragraph{9. What is left behind.} A run writes a directory
\texttt{results/run\_\textless epoch\textgreater\_\textless dataset
hash\textgreater/} holding: \texttt{config.json}, the arm configuration
actually used---ranking function, tier mode, own-post exclusion, corpus
snapshot statistics, agent-rules hash---read back from the objects the run
built rather than from the environment it was launched with;
\texttt{per\_query.jsonl}, one row per question per arm with the ranked
document ids, every metric, latency, model calls and tokens;
\texttt{qrels.json}, the labels used by \emph{this} run alone, so its numbers
recompute from its own artifacts without the accumulated cache;
\texttt{aggregate.csv} and \texttt{report.md}. Each table in
Section~\ref{sec:results} names the run id it was computed from. Two senses of
recovery apply to it, and they are not the same. Every metric recomputes from
that directory alone, for any run: the ranked document ids and the judgments
are both on disk, so a reader can re-derive the numbers, re-score them under a
different metric, or pool them with another run's judgments without touching
the endpoints. Re-executing the pipeline and arriving at the same ranked lists
is the stronger sense, and it holds for runs made in replay, not for runs made
in record mode against the live feeds.



\section*{S2\quad Arm Comparability}


Set side by side, the two arms do not hold the same capabilities. The
multi-agent arm rewrites the query, fetches both phrasings, scores its own
retrieval, and---in the full architecture---retries and accumulates memory; the
baseline does none of these. A reader is entitled to ask whether the two arms
are performing the same task at all, and if they are not, what a difference
between them measures.

The asymmetries are of two kinds, and only one of them is a defect. Query
rewriting, union fetch, and adaptive retry \emph{are} the treatment: they are
the independent variable, and removing them from the multi-agent arm would
leave nothing to evaluate. What makes their comparison sound is not symmetry of
capability but symmetry of everything else. Both arms retrieve from the same
endpoints through the same Retriever, cut to the same $k$, synthesize through
the same prompt with the same model, are scored by the same judge on the same
questions, and are compared pairwise within a single run so that corpus drift
cancels. The baseline is an ablation of the coordination layer, not a different
system.

The second kind of asymmetry is not part of the treatment and does leak into
measurement: the bespoke retrieval score of Section~\ref{sec:two-metrics} and
the answer rendering of Section~\ref{sec:scaled-eval}. The bespoke score is the
Evaluator's internal control signal---the same quantity that gates the
retry---and the baseline has no counterpart to it, so reporting it as the
comparison metric grades the treatment with its own ruler; Section~\ref{sec:metric-threat}
retires it for that reason. Answer rendering is neutralized rather than
defended, by the \texttt{marag\_llm} arm. Confounds of this kind must be
removed; treatment asymmetries must not be.

Removing them one at a time is what the ablation ladder in
Table~\ref{tab:ladder} does. Each rung differs from the one below it by exactly
one capability, so each adjacent pair is a single-factor comparison rather than
a comparison of two architectures. Rung A3 additionally runs the Evaluator, but
in a prose arm the Evaluator's output is only its quality score: it changes
neither the retrieved documents nor the answer, so A2 versus A3 isolates the
union fetch and nothing else.



The same reasoning applies to rendering, which is not a property of the
multi-agent arm but a factor crossing it. Of the four cells of
rendering~$\times$~retrieval, three were populated in
Section~\ref{sec:scaled-eval}: baseline retrieval with prose
(\texttt{single\_agent}), multi-agent retrieval with prose
(\texttt{marag\_llm}), and multi-agent retrieval with a template
(\texttt{marag}). The format effect reported there is therefore an inference
across two arms rather than an estimated main effect. The fourth cell---baseline
retrieval rendered through the Evaluator's template---closes the design; it
requires no generation call, since the template is assembled deterministically
from the retrieved documents, and it is available in the released harness as
\texttt{single\_agent\_template}. With all four cells, rendering and retrieval
become main effects with an interaction term.

A2, A4, and the fourth cell are now measured, together with A1g, and
Section~\ref{sec:ladder-results} reports them. The run satisfies the constraint
this design imposes: because the retrieval endpoints are live, every arm
executed against a recorded corpus snapshot in replay, so the rungs differ by
capability and not by corpus drift. One constraint remains open. If the
Self-Improvement Memory is ever attached as a further rung, term-weight
learning makes each question's outcome depend on the questions preceding it;
per-question differences then cease to be exchangeable and the paired tests
used throughout this section no longer apply as written.



\section*{S3\quad Evaluation Methodology for Retrieval-Augmented Systems}

Retrieval evaluation inherits from TREC the use of pooled relevance judgments:
candidates are collected from the runs of all participating systems and judged
once, so no system is scored against a judgment set built around its own output.
Pooling is nonetheless incomplete. Zobel found that TREC measured relative system
performance reliably but overestimated recall, because many relevant documents
had likely never been found~\cite{zobel1998reliable}, and Buckley and Voorhees
showed that standard measures are not robust to substantially incomplete
judgments~\cite{buckley2004incomplete}. Both concerns apply to a comparison
between pipelines that fetch different document sets from live endpoints, and
they are why the judgments used here are pooled across both systems and all
ablation arms.

Benchmarks supply the other half of the practice. BEIR measures zero-shot
transfer across heterogeneous retrieval tasks~\cite{thakur2021beir}; KILT fixes a
single Wikipedia snapshot and requires provenance~\cite{petroni2021kilt}; the
CRAG benchmark spans five domains and eight question categories, with mock APIs
simulating web and knowledge-graph search~\cite{yang2024}; and FreshQA targets
questions whose answers change over time, the closest published analogue to
software update questions~\cite{vu2024freshllms}. Reference-free scorers such as
RAGAS~\cite{es2024ragas} and ARES~\cite{saadfalcon2024ares} estimate
faithfulness and relevance without gold answers.

Two aspects of that practice bear on the metric inversion reported here. The
first is judge independence, and the IR literature has by now measured it
precisely enough to be quantitative about what an LLM judge can and cannot
support. The encouraging result is at the level of system \emph{ordering}:
UMBRELA, an open-source reproduction of Bing's assessor, induces rankings that
agree with NIST qrels on the TREC Deep Learning tracks at Kendall's $\tau$ of
0.873 to 0.944~\cite{upadhyay2024umbrela}, and Thomas et al.\ report run-level
$\tau$ of 0.77 to 0.86 over 110 TREC-Robust runs~\cite{thomas2024preferences}.
The discouraging result is at the level of individual \emph{labels}, which is
what a qrel file actually contains: UMBRELA's per-item Cohen's
$\kappa$~\cite{cohen1960coefficient} against
human assessors is only 0.31 to 0.37 on the four-point
scale~\cite{upadhyay2024umbrela}, and Faggioli et al.\ measure $\kappa=0.26$ on
the TREC-DL 2021 pool against $\kappa=0.52$ between two groups of human
assessors~\cite{faggioli2023perspectives}. Agreement is also asymmetric:
LLM judges over-label relevance, so a machine ``not relevant'' label is
considerably more trustworthy than a machine ``relevant'' one, and they can be
induced to grade a passage of randomly sampled words as perfectly relevant
simply by injecting the query string into
it~\cite{alaofi2024fooled}. Faggioli et al.\ accordingly place practice on a
spectrum from human judgment through AI assistance and human verification to
full automation, and decline to endorse the fully automated
endpoint~\cite{faggioli2023perspectives}; Soboroff argues the stronger position
that whatever generates the qrels also defines the ceiling, so a system better
than the judge cannot be measured as
such~\cite{soboroff2025dontuse}. He also names the condition under which that
objection lapses, and it is the one pooling supplies: judgments gathered from
many systems carry more relevance information than any single system, so ``we
can measure any individual systems with less information about relevance than
the collective pool''~\cite{soboroff2025dontuse}. The concession is guarded ---
he requires the privilege to be explicit, a judge fine-tuned on manual
judgments for the evaluation topics being his example, and warns against
assuming it from training data one cannot inspect --- but pooling across
systems is the form of it obtainable without human assessors, which is why
every judgment reported here is pooled rather than collected per system.

Circularity is the specific form of this problem that a multi-agent pipeline
invites, and its cost has been simulated. Clarke and Dietz re-ranked all 75 runs
of the TREC 2024 RAG track with UMBRELA and then scored them with UMBRELA:
agreement with human judgments fell from $\tau \approx 0.84$ to 0.38 at rank~10
and to $-0.40$ at rank~5, twelve systems exceeded nDCG 0.95 where manual
assessment placed them at 0.68 to 0.72, and a deliberately constructed run rose
from 28th under human assessment to 5th under automatic
assessment~\cite{clarke2024cannot}. The same effect is documented at the level
of preference: judges recognize and favor text they themselves
generated~\cite{panickssery2024self}, and LLM-as-judge studies report position,
verbosity, and self-enhancement biases~\cite{zheng2023judge}. A judge that is
also the pipeline's generator, or that is prompted with the pipeline's own
notion of relevance, is not independently measuring it. There is one direction
in which this literature works in favour of a finding rather than against it:
because LLM judges are documented to penalize lexically grounded retrieval
relative to LLM-mediated retrieval~\cite{faggioli2023perspectives}, a result in
which the plainer, non-rewriting system \emph{wins} is not the result this bias
would manufacture. The second
is abstention. The CRAG benchmark scores an accurate answer $1$, a missing answer
$0$, and an incorrect answer $-1$, explicitly preferring an admission of
ignorance to a confident
error~\cite{yang2024}; unanswerable-question benchmarks draw the same
distinction on the reading-comprehension side~\cite{rajpurkar2018squad2}. A
keyword-overlap score registers neither distinction: it rewards the presence of
question terms in the retrieved text whether or not that text answers the
question, and assigns no cost to a fluent answer assembled from documents that do
not support one.

The general caution is not new. Armstrong et al.\ analysed a decade of reported
ad-hoc retrieval improvements and found little evidence of cumulative progress,
because baselines were generally weak --- often below the median original TREC
system --- and questioned the value of a statistically significant gain over such
a baseline~\cite{armstrong2009improvements}; Lipton and Steinhardt catalogue the
reporting patterns that let such results stand~\cite{lipton2019trends}. A
retrieval-quality score defined by the authors of the system it evaluates is the
limiting case, having neither external calibration nor a cross-system pool. What
this paper adds is a worked instance in which such a score and a standard one
disagree by roughly a factor of five in opposite directions over the same runs,
with the document-level audit that establishes which was right.



\section*{S4\quad The Multi-Ecosystem Benchmark}

The evaluation in the conference version used 50 questions drawn from a single
community source, which is too few to separate an architectural effect from
noise and too narrow to say anything about ecosystems other than the ones those
questions happened to mention. We therefore constructed a larger benchmark, and
we describe its construction here in enough detail to expose its biases, because
a benchmark assembled by the authors of the system it evaluates warrants that
scrutiny.

\subsubsection{Construction}

The benchmark contains 500 questions arranged as a $24 \times 5$ grid of
software ecosystem by question category, filled to 20 or 21 questions per
ecosystem and exactly 100 per category. It was built in two stages from one
cache and one seed --- 300 questions first, then 500 --- and the selection is
greedy over deterministically ordered pools, so the 300 are a strict subset of
the 500; where a measurement below is reported at $n=300$, it is on that
subset. The 24 ecosystems span mobile and desktop
platforms (Apple iOS, Apple macOS, Android, Windows), four Linux distributions
(Ubuntu, Debian, Fedora, Arch) and the kernel itself, browsers (Firefox,
Chrome), container infrastructure (Docker, Kubernetes), language package
ecosystems (npm/Node, pip/Python, cargo/Rust), self-hosted and home applications
(Home Assistant, NAS/self-hosted, WordPress), databases (PostgreSQL, MySQL,
NoSQL), developer tooling, and LLM/AI model releases. The five categories are
releases, bugs, security, community, and general.

Every question is mined from live data rather than authored: 249 are verbatim
titles of real community posts, and 51 are generated from real release or
advisory records. No question is a template with a hand-written stem. The
builder draws from a pool of 9,923 community posts and the release and advisory
endpoints described in Section~\ref{sec:datasources}, and it is deterministic:
a fixed seed, an on-disk response cache, and an \texttt{--offline} mode that
replays the cache, so two runs of the builder produce byte-identical output.
Each record carries the endpoint and record type it came from, so mined and
generated questions remain distinguishable in any analysis.

\subsubsection{Two Upstream Data Defects}

Two properties of the upstream feeds would have injected false questions into
the benchmark had they been consumed naively. We report them because they are
properties of the public data sources rather than of our pipeline, and anyone
building on the same feeds will encounter them.

This is not an isolated quirk of one feed. Vulnerability-database metadata is
known to be noisy: affected-version and severity records in the NVD contain
substantial inaccuracy~\cite{croft2023dataquality,croft2022severity,anwar2022nvd},
and inconsistencies between advisory text and structured metadata have been
documented at scale~\cite{dong2019inconsistencies}. Downstream models trained on
such labels inherit the noise~\cite{chakraborty2022vulndetection}. What follows is
the form the problem takes in the feeds used here.

First, on advisory records the product field names the \emph{index token} under
which the advisory was filed, not the affected product. An advisory for ``Mocha
Telnet Lite for iOS 4.2'' is stored with product \texttt{iOS} and version
\texttt{4.2.0}; advisories concerning Velero are stored under
\texttt{Kubernetes}. Templating those fields directly yields questions such as
\emph{``Is iOS v4.2.0 vulnerable?''}, which assert a falsehood in the question
itself and would be scored against a system for failing to confirm it. The
builder instead parses the affected product and version out of the advisory
text, keeps a record only when the parsed product belongs to a tracked
ecosystem, and discards the rest: of 3,113 advisory rows, 548 parsed and 268
could be confidently attributed.

Second, a minority of records carry malformed dates (for example
\texttt{23.30.13102017} and \texttt{2025[29]0213}). Since date agreement is one
of the two decisive facts in this domain, these records are dropped rather than
normalized by guessing.

We also stopped trusting the feed's own topical flags. Its CVE-relatedness flag
fires on posts such as \emph{``Best budget AI subscription / API tokens?''}, so
it is used only to break ties within the security category rather than to assign
one.

\subsubsection{Ground Truth and Its Limits}

Of the 500 questions, 72 carry a reference answer (51 of the first 300), and
only those generated from release or advisory records do --- a community post
title states a question but does not come with an authoritative answer.
Ground-truth coverage is therefore concentrated in the releases and security
categories, and any correctness or hallucination measurement over this
benchmark is computed on that subset rather than on all 500 questions.

A community post title comes with something else: the URL of the post it was
taken from, which the live community feed the system retrieves from will return
as a document. Section~\ref{sec:leak} reports what that did to every retrieval
measurement and how the harness now excludes it; the benchmark file carries the
URL on every record so that the exclusion, and the post-hoc correction of runs
made without it, are reproducible.

This has a consequence for sampling that we did not anticipate and that a
reader reusing the benchmark should know. A subset selected to balance ecosystem
and category, without regard to ground truth, can leave almost none: our first
100-question stratified selection carried a reference answer on 2 questions. A
selection that prefers ground-truth records within each grid cell retains 28.
Reported retrieval metrics are unaffected --- they rest on pooled relevance
judgments, not reference answers --- but any answer-correctness figure requires
the second kind of selection, and we report retrieval and answer-faithfulness
results rather than correctness for this reason.

\subsubsection{Coverage Gaps}

One of the 120 grid cells (Arch $\times$ security) had no live candidate
available at build time. Rather than synthesize one, the builder logs the gap
and fills the category quota from other ecosystems, so the category totals hold
exactly while that single cell is empty. Nineteen product-name lookups returned
404 because the product is not in the vendor catalogue (npm, pip, cargo,
Synology, Home Assistant, NixOS, GitLab among them); this is a catalogue
coverage limitation rather than an endpoint failure, and the affected ecosystems
are represented through their community and release records instead.


\section*{S5\quad Software Ecosystem Data Sources}

The system retrieves information from software ecosystem documents
obtained through multiple API endpoints.

\textbf{Verified Sources}
\begin{itemize}
\item \texttt{/\{vendor\}}: vendor registry (6,578 entries)
\item \texttt{/v/}: software release notes (31,958 entries)
\item \texttt{/cve}: vulnerability advisories (24,139 entries)
\end{itemize}

\textbf{Community Sources}
\begin{itemize}
\item \texttt{/reddit}: community discussions (208,466 posts)
\item \texttt{/positive}: software update risk discussions (2,637 posts)
\item \texttt{/by-subreddit?q=\{vendor\}}: vendor-specific discussions
\item \texttt{/questions}: software-related user questions
\end{itemize}

To support vendor-aware retrieval, a registry of software vendors and
associated subreddits is loaded at startup and cached locally.

The entry counts above are those measured when the conference version was
written. They are not re-measurable: the endpoints return query results
without a collection total, so a reader cannot check them against the live
service. The one count that can be checked has moved substantially ---
\texttt{/api/c/names}, the product vocabulary the vendor matcher is built
from, returned 14,223 names then and returns 1,348 today (2026-09-23, full
response, not paginated). We report the original figures rather than silently
restate them, and note that any absolute corpus size in this section should be
read as a measurement of the service at that time rather than a stable
property of it. Retrieval results in this paper are unaffected: every run
either records its own corpus snapshot or is replayed from one
(Section~\ref{sec:tier-control}).



\section*{S6--S9\quad Superseded Measurements at $n=100$}

These tables report the 100-question measurements the paper's later runs at
$n=500$ and $n=1{,}000$ supersede. They are retained because the main text
cites them when tracing how each conclusion was reached.

\begin{table}[t]
\caption{Ranking ablation repeated on the repaired pipeline, $n=100$, against a
frozen corpus snapshot replayed identically for all three arms
(\texttt{corpus\_misses}~$=0$, \texttt{frozen}~$=$~true). Unlike
Table~\ref{tab:ranking-ablation} these rows are strictly paired. With the union
fetch in place, stronger ranking now moves nDCG@3 and MRR \emph{together} ---
the divergence that localized the defect is absent, which is what the diagnosis
predicts. Levels in this table are inflated by the own-post leak of
Section~\ref{sec:leak} (72--73\% of top-$k$ lists on this frozen corpus); the
within-table contrasts, which compare arms that leak at the same rate, are what
it is cited for.}
\label{tab:scaled-ablation}
\begin{tabular}{llrrrr}
\toprule
Ranking function & System & nDCG@3 & nDCG@5 & MRR & Recall@5 \\
\midrule
None            & multi-agent  & 0.486 & 0.473 & 0.618 & 0.341 \\
                & single-agent & 0.507 & 0.493 & 0.646 & 0.371 \\
Okapi BM25      & multi-agent  & 0.627 & 0.596 & 0.793 & 0.387 \\
                & single-agent & 0.631 & 0.606 & 0.806 & 0.396 \\
Embedding cosine & multi-agent  & \textbf{0.656} & \textbf{0.615} & \textbf{0.817} & 0.387 \\
                 & single-agent & 0.646 & 0.612 & 0.808 & 0.388 \\
\bottomrule
\end{tabular}
\end{table}

\begin{table}[t]
\caption{Answer quality and cost ($n=100$), LLM-judged, mean $\pm$ 95\% CI.
\texttt{marag} and \texttt{marag\_llm} share a retrieval pipeline and differ
only in how the answer is produced. Quality is measured on the frozen-corpus
replay; latency is the per-question median from the live recording pass that
produced that snapshot, since a replay reads its documents from cache and its
timings measure nothing.}
\label{tab:scaled-answer}
\begin{tabular}{lrrr}
\toprule
System & Faithfulness & Answer relevance & Latency (s) \\
\midrule
marag (template answer)   & 0.748 $\pm$ 0.036 & 0.995 $\pm$ 0.004 & 18.8 \\
marag\_llm (prose answer) & 0.919 $\pm$ 0.035 & 0.968 $\pm$ 0.027 & 24.7 \\
single\_agent             & 0.929 $\pm$ 0.032 & 0.973 $\pm$ 0.021 & 8.4 \\
\bottomrule
\end{tabular}
\end{table}

\begin{table}[t]
\caption{The same three systems on a second 100-question selection, one that
retains 28 reference answers, under the same frozen-corpus protocol (6{,}227
snapshot hits, no miss on a document host). Values are mean $\pm$ 95\% CI;
correctness is computed only where a reference answer exists, with $n$ given
per arm; latency is the median from the live recording pass. The nominal
ordering on retrieval is the reverse of Table~\ref{tab:scaled-retrieval}, and
remains so after the repair --- which is the point of reporting this sample.
Retrieval levels carry the own-post leak (72--73\% of top-$k$ lists; leak-free
nDCG@3 0.330 / 0.330 / 0.333).}
\label{tab:gt-run}
\begin{tabular}{lrrrr}
\toprule
System & nDCG@3 & Faithfulness & Correctness & Latency (s) \\
\midrule
marag         & 0.795 $\pm$ 0.071 & 0.735 $\pm$ 0.037 & 0.324 $\pm$ 0.109 ($n=27$) & 16.2 \\
marag\_llm    & 0.795 $\pm$ 0.071 & 0.939 $\pm$ 0.030 & 0.392 $\pm$ 0.160 ($n=25$) & 16.5 \\
single\_agent & 0.773 $\pm$ 0.075 & 0.932 $\pm$ 0.031 & 0.363 $\pm$ 0.134 ($n=27$) & 8.2 \\
\bottomrule
\end{tabular}
\end{table}

\begin{table}[t]
\caption{Paired per-question differences, Wilcoxon signed-rank. ``Better'' and
``worse'' count the questions on which the second arm beat or lost to the
first; the remainder tie. Rows marked $\dagger$ report a $p$ below the smallest
value the test can attain at that number of non-zero pairs (0.0625 at five
pairs, 0.25 at three): they are null, not significant.}
\label{tab:ladder-paired}
\begin{tabular}{llrrrl}
\toprule
Comparison & Metric & Mean $\Delta$ & Better & Worse & $p$ \\
\midrule
A1 $\to$ A1g  & nDCG@3   & $+0.049$ & 14 & 5  & \textbf{0.0070} \\
A1 $\to$ A1g  & Recall@3 & $+0.051$ & 12 & 4  & \textbf{0.0035} \\
A1 $\to$ A1g  & nDCG@5   & $+0.042$ & 15 & 9  & \textbf{0.0268} \\
A1 $\to$ A1g  & Recall@5 & $+0.037$ & 13 & 9  & 0.0534 \\
A1 $\to$ A1g  & MRR      & $+0.023$ & 5  & 0  & 0.0431$^\dagger$ \\
A1 $\to$ A1g  & nDCG@1   & $+0.027$ & 3  & 0  & 0.1088$^\dagger$ \\
\midrule
A1 $\to$ A3   & nDCG@5   & $+0.018$ & 9  & 10 & 0.8721 \\
A1 $\to$ A3   & Recall@5 & $+0.007$ & 6  & 8  & 1.0000 \\
A1 $\to$ A3   & MRR      & $+0.016$ & 4  & 4  & 0.4838 \\
\midrule
A1g vs.\ A3   & nDCG@5   & $-0.024$ & 12 & 20 & 0.2172 \\
\bottomrule
\end{tabular}
\end{table}

\section*{S10\quad Multi-Agent RAG Frameworks, Reproducibility Toolkits, and the Self-Reflective Arm}

Multi-agent RAG systems distribute retrieval and reasoning across
specialized agents. MA-RAG employs planner, extraction, and question
answering agents~\cite{nguyen2025marag}. MAIN-RAG separates document
selection and answer generation through dedicated predictor and judge
agents~\cite{wang2024mainrag}. More generally, frameworks such as
AutoGen~\cite{wu2023autogen} and smolagents~\cite{huggingface2024agents}
provide infrastructure for agent orchestration and tool use.

Four recent systems are close enough to the pipeline studied here to warrant
individual comparison. MA-RAG chains Planner, Step Definer, Extractor, and QA
agents that hand chain-of-thought state to one another, and reports that
LLaMA3-8B under this arrangement outperforms larger standalone models on
multi-hop benchmarks~\cite{nguyen2025marag}; it shares our staged specialists
and 8B-class local model but is evaluated on Wikipedia-derived corpora with no
memory across questions. RAGentA feeds hybrid sparse--dense retrieval into a
draft-answer agent, a triplet filter, a citing answer agent, and a completeness
checker that reformulates the query when evidence is
lacking~\cite{besrour2025ragenta}; its filter-and-cite stages correspond to our
reranker and source-traced answers, but its reformulation runs \emph{after} a
first retrieval over the original wording, so the original query is always
searched --- the property whose absence produced the defect reported in
Section~\ref{sec:rewrite-roles}. Amber maintains an agent-updated memory that steers
an adaptive information collector and a multi-granular filter across retrieval
iterations~\cite{qin2025amber}; it is the nearest analogue to our
self-improvement memory, though Amber's memory lives inside a single question's
loop whereas ours persists across successive queries and is updated from their
outcomes. Shaikh ablates decomposition, adaptive routing, loop depth, and
cross-encoder reranking on a local Qwen2.5-7B and finds fixed hybrid retrieval
beating adaptive routing, with decomposition and reranking contributing only
modest gains~\cite{shaikh2026dissecting} --- a design-space result that agrees
with our own finding that decomposition per se did not produce the originally
reported improvement, though our result arrived as a defect diagnosis rather
than an ablation. Finally, Mishra et al.\ systematise agentic RAG as a
partially observable decision process along planning, retrieval orchestration,
memory, and tool-use axes, and catalogue retrieval misalignment and memory
poisoning as systemic risks~\cite{mishra2026sok}; the replacing rewriter
described in this paper is a concrete instance of the former.

Most prior evaluations focus on general question-answering benchmarks, whereas
this work evaluates a multi-agent architecture for software update question
answering over a live, heterogeneous corpus.

\subsection{Multi-Agent RAG and Reproducibility Infrastructure}

The self-correcting designs reviewed above --- Self-RAG's reflection
tokens~\cite{asai2024selfrag}, CRAG's lightweight retrieval
evaluator~\cite{yan2024crag}, Adaptive-RAG's complexity-based
routing~\cite{jeong2024adaptiverag}, and FLARE's confidence-triggered retrieval
during generation~\cite{jiang2023flare} --- decide \emph{whether} and \emph{how
often} to retrieve, and how to filter what returns. None of them, as published,
checks whether the text handed to the retriever preserves the user's own
phrasing.

Two toolkits now make much of this work runnable. FlashRAG is an MIT-licensed
Python toolkit that reimplements RAG algorithms under one interface --- 16
methods over 38 datasets as published~\cite{jin2025flashrag}, 23 algorithms over
36 in the current release, including Self-RAG, Adaptive-RAG, and FLARE
pipelines~\cite{flashrag2025repo}. RAGLAB reproduces six algorithms across ten
benchmark datasets~\cite{zhang2024raglab} and publishes fine-tuned
\texttt{selfrag\_llama3-8B} checkpoints with VLLM-based and quantized generation
configurations~\cite{raglab2024repo}.

We did not run a head-to-head comparison against these baselines, for a reason
that is a property of our setting rather than a criticism of the toolkits: each
assumes a static corpus embedded and indexed in advance. Self-RAG's released
retrieval component is Contriever~\cite{izacard2022contriever} over the DPR
Wikipedia passage split~\cite{karpukhin2020dpr} with precomputed embeddings; its
authors note that retrieval over full English Wikipedia requires large memory
and multiple GPUs, and report testing their on-demand retrieval script on a
configuration with 100\,GB of RAM, while observing that less should
suffice~\cite{selfrag2023repo}. RAGLAB's ColBERT retrieval server is documented
as requiring at least 60\,GB of RAM~\cite{raglab2024repo}. CRAG's retrieval
evaluator is a fine-tuned T5-large model, and its confidence thresholds are set
separately for each evaluation dataset --- $(0.59, -0.99)$ for PopQA and
$(0.95, -0.91)$ for biography generation, among others~\cite{yan2024crag} --- so
transferring the method to a new domain requires re-tuning against labelled data
for that domain. The system studied here retrieves from live release, advisory,
and community APIs whose contents change between runs and for which no pre-built
index exists: there is no static split to embed, and per-dataset thresholds have
no analogue. The comparison we can make is against a single-agent baseline
issuing the same calls to the same endpoints in the same time window, which
isolates the architectural variable at the cost of not situating either system
against published leaderboard numbers. We record this as a limitation rather than
resolving it.

One partial step toward resolving it is released with the harness. Since the
obstacle is the retrieval stack rather than the methods, we reimplemented the
inference-time mechanisms both papers introduce --- Self-RAG's per-passage
relevance and support critiques, and CRAG's three-way corrective control flow
--- over our own retrieval layer, with the synthesis model held constant and
CRAG's web-search branch replaced by declining, since no web fallback exists in
a fixed-corpus comparison. It is available as an arm of the evaluation harness
(\texttt{--generators selfreflective}). We report no results for it: its
evaluation was not complete at submission, and reporting a partial run of a
baseline would be worse than reporting none. What the arm establishes now is
only that the mechanisms can be evaluated in this setting once the retrieval
obstacle is set aside, and that a future comparison need not port a Wikipedia
index to make it. It is a reimplementation and not a reproduction: neither
published checkpoint is involved, so it can never settle how those systems
perform, only how their mechanisms behave over these sources.

Completing that evaluation carries a condition worth stating in advance. In the
configuration released here the arm's critiques run the same model that judges
answers for scoring, and the critique loop decides both which documents the
answer is drafted from and whether to answer at all. A judge scoring answers it
helped shape is the self-preference case documented for LLM
evaluators~\cite{panickssery2024self,zheng2023judge}, so the arm's
\emph{answer-quality} numbers under such a judge would not be independent of it.
Whoever finishes the run should score it with a model independent of the critic.

The arm's \emph{retrieval} metrics are not exposed to this, by a deliberate
choice in the harness worth recording because the opposite is the natural
implementation: the arm reports the full retrieved list rather than the
critique-filtered one, so nDCG and recall are computed over what it retrieved
and not over what its critic approved. Scoring only the kept documents would
flatter precision by hiding the discards, and would additionally make the
retrieval numbers circular in the sense Clarke and Dietz
measure~\cite{clarke2024cannot}. The harness records each critique decision, so
the number of documents the critic withheld from drafting is recoverable per
query; because a hedged or negated critique keeps the document
(\texttt{SelfReflectiveRAG.is\_relevant} discards only on a non-negated
\textsc{irrelevant} verdict, via \texttt{find\_verdict}), that count is a lower
bound on the critic's filtering.


\end{document}
